\chapter{Kathará}
\section{Introduzione}
\begin{defn}[Kathará]
  Kathará è un framework open source basato su container Docker che permette di creare, configurare ed eseguire reti virtuali,
  integrando in un unico ambiente tecnologie come \textbf{SDN} (Software Defined Network) e \textbf{NFV} (Network Function Virtualization) \cite{8406267}.
\end{defn}
Ogni nodo di rete è simulato con un container. Il comportamento del nodo dipende dall'immagine usata,
permettendo l'uso di protocolli tradizionali (OSPF, BGP, ecc), SDN o servizi applicativi (DNS, Apache, ecc).
I vari dispositivi comunicano tra loro attraverso LAN virtuali di livello 2, riproducendo il comportamente di un collegamento fisico.\\
\vspace*{0.01\textheight}
La scelta di basarsi su container invece che su macchine virtuali, rende Kathará uno strumento molto leggero.
Ogni container richiede risorse nettamente inferiori rispetto a una macchina virtuale tradizionale, permettendo di simulare topologie complesse su hardware meno performante.
Inoltre, permette di sfruttare bene tecnologie come NFV (Network Function Virtualization), che utilizza funzioni software per sositutire dispositivi hardware,
e SDN (Software Defined Network), che permette di separare la logica di controllo della rete, dalla logica di invio dei pacchetti.

\section{Architettura di Kathará}
L'architettura di Kathará è composta da tre moduli principali: \textbf{Scripts}, \textbf{Operations} e il \textbf{Container Platform Adapter} \cite{8406267}.
\begin{figure}[!htb]
    \centering
    \includegraphics[width=0.65\linewidth]{Images/kathara/Kathara.png}
    \caption{Architettura di Kathará \cite{8406267}}
\end{figure}
\newpage
Il modulo Scripts contiene le operazioni di base che utente può fare per avviare, fermare e mettere in pausa i nodi di rete. Si dividono in due tipologie:
\begin{enumerate}
  \item \textbf{ltools:} valgono per tutti i nodi inseriti nel file di configurazione (\texttt{lstart}, \texttt{lclean}).
  \item \textbf{vtools:} valgono per un singolo nodo (\texttt{vstart}, \texttt{vclean}).
\end{enumerate}
Il modulo Operations, invece, racchiude alcune funzionalità che servono a Kathará per funzionare, come:
\begin{itemize}
  \item \textbf{Parser:} traduce i file di configurazione o i parametri passati da riga di comando in nodi della rete con eventuali collegamenti. 
    Si occupa anche di tradurre ulteriori opzioni, come risorse assegnate a ogni nodo o immagini Docker da utilizzare.
  \item \textbf{Start-Up Operations:} serve a garantire compabilità fra i diversi sistemi.
  \item \textbf{Common Libraries:} genera una serie di operazioni indipendenti da eseguire.
\end{itemize}
Infine, il modulo del Container Platform Adapter converte i comandi generati in precedenza in direttive Docker (o altre piattoforme) per creare e collegare fra di loro i container.
L'output dell'adapter ritorna al modulo Scripts, che lo invia al Docker daemon. \\
\vspace*{0.01\textheight}
Nei sistemi Unix, per ragioni di sicurezza, i comandi devono prima essere approvati da un wrapper che scarta qualsiasi comando non ritenuto sicuro.

\section{Configurazione}
Supponiamo di avere un rete con tre nodi: due host e uno switch. Tramite Kathará è possibile crearli e avviarli tramite riga di comando, con i vtools, oppure, tramite file di configurazione, con gli ltools \cite{8406267}.
Questa seconda modilità è quella che userò per configurare la rete.\\
\vspace*{0.01\textheight}
Il primo file da creare è quello in cui vengono definiti i nodi e se necessario altre opzioni, come le immagini dei container Docker. 
\begin{kathara}[lab.conf]
  # Collegamento fra un host e porte dello swicth
  pc1[0]="A"
  sw1[0]="A"

  pc2[0]="B"
  sw1[1]="B"

  # Immagine per lo switch
  sw1[image]="kathara/quagga"
\end{kathara}
Successivamente, è possibile creare a avviare i container con il comando \texttt{kathara lstart}, partendo dalle immagini di default o da quelle specificate. 
Kathará, inoltre, aprirà per ogni nodo un terminale, in modo da poter fare modifiche o test. \\
\vspace*{0.01\textheight}
In alternativa, è possibile creare dei file di configurazione per ogni nodo, scrivendo al suo interno una lista di comandi da far fare in automatico all'accensione del container.
Per l'host \texttt{pc1}, nel suo file di configurazione si può assegnare un ip e accendere la rispettiva interfaccia.
\begin{kathara}[pc1.startup]
  # Assegnare un ip
  ip addr add 192.168.1.1/24 dev eth0

  # Accendere l'interfaccia eth0
  ip link set dev eth0 up
\end{kathara}
Tramite questa modalità è possibile configurare pure il secondo host e lo swicth. Una volta digitato il comando \texttt{kathara lstart}, tutti i file di configurazione verrano usati per creare la rete.

\subsection{Riproduzione dello scenario HTB in Kathará}
Sfruttando Kathará, si possono riproporre gli esempi del capitolo precedente sull'utlizzo delle qdisc tramite il comando \texttt{tc}, testandolo con diversi nodi.
In particolare, lo scenario su HTB è utile per dimostrare che è possibile fare traffic shaping con questo framework.
\vspace*{0.01\textheight}
Facendo riferimento allo scenario dell'esercizio \ref{sec:htb}, il primo file da creare sarà \texttt{lab.conf}, in cui dovranno essere definiti i tre clienti e il server con i servizi.
\begin{kathara}[lab.conf]
  # Clienti
  c1[0] = "A"
  c2[0] = "A"
  c3[0] = "A"
  
  # Server
  server[0] = "A"
\end{kathara}
Specificando per ognuno il dominio \texttt{"A"}, Kathará simulerà il comportamente di una rete locale con uno swicth che collega tutti i nodi.
Successivamente, per ogni host, dovrà essere creato un file che conterrà tutti i comandi necessari per poter comunicare fra di loro e nel caso del server pure quelli utili per il traffci shaping.
\begin{kathara}[c1.startup]
  # Assegnazione ip
  ip addr add 10.0.0.1/24 dev eth0

  # Accensione interfaccia eth0
  ip link set dev eth0 up
\end{kathara}
Questa configurazione sarà uguale anche per gli altri clienti, ad eccezzione dell'indirizzo ip differente. Invece, in \texttt{server.startup}, verranno inseriti pure i comandi con \texttt{tc}.
\begin{bash}[server.startup]
  # Assegnazione ip
  ip addr add 10.0.0.10/24 dev eth0

  # Accensione interfaccia eth0
  ip link set dev eth0 up

  # Nodo root e genitore
  tc qdisc add dev eth0 root handle 1: htb default 30
  tc class add dev eth0 parent 1: classid 1:1 htb rate 100mbit ceil 100mbit

  # Classe sito web
  tc class add dev eth0 parent 1:1 classid 1:10 htb rate 50mbit ceil 100mbit

  # Classe servizio streaming
  tc class add dev eth0 parent 1:1 classid 1:20 htb rate 35mbit ceil 60mbit

  # Classe servizio file hosting
  tc class add dev eth0 parent 1:1 classid 1:30 htb rate 15mbit ceil 20mbit
\end{bash}
Infine, all'interno di quest'ultimo file, o sul container una volta avviato, dovranno essere creati pure i filtri per classificare i flussi.
\begin{bash}[server.startup]
  # Filtri per sito web
  tc filter add dev eth0 protocol ip parent 1:0 prio 1 u32 match ip sport 80 0xffff flowid 1:10
  tc filter add dev eth0 protocol ip parent 1:0 prio 1 u32 match ip sport 443 0xffff flowid 1:10

  # Filtro per servizio streaming
  tc filter add dev eth0 protocol ip parent 1:0 prio 2 u32 match ip sport 554 0xffff flowid 1:20

  # Filtro per servizio file hosting
  tc filter add dev eth0 protocol ip parent 1:0 prio 3 u32 match ip sport 21 0xffff flowid 1:30
\end{bash}

\subsubsection{Risultati}
Per capire gli effetti che il traffic shaping ha portato alla rete, è utile confrontare i tempi di invio di dati con o senza HTB. La formula per calcolare i tempi di invio ideali è:
\[
  t = \frac{\text{dim} \times 8}{\text{rate}}
\]
Sul nodo del server è stato creato un file di 100 MB e poi è stato posto in ascolto sulle tre porte: \texttt{443}, \texttt{554} e \texttt{21}. 
Ogni client ha scaricato contemporaneamente il file, tramite il comando \texttt{nc}, collegandosi alla porta del proprio servizio. \\ 
Per ciascun trasferimento è stato misurato il tempo reale.
\begin{center}
  \begin{tabular}{ |c|c|c| }
    \hline
    \textbf{Classe} & \textbf{Tempo Ideale} & \textbf{Tempo Reale} \\
    \hline
    1:10 & $\sim$ $16s$ & $19.20s$ \\
    \hline
    1:20 & $\sim$ $23s$ & $23.33s$ \\
    \hline
    1:30 & $\sim$ $53s$ & $48.47s$ \\
    \hline
  \end{tabular}
\end{center}
Nell'esempio, si può notare come la classe \texttt{1:30} abbia impiegato meno tempo di quello ideale. Questo è dovuto al fatto che,
dopo che il primo cliente ha finito di leggere i \texttt{100MB} di file, le altre due classi si sono divise la banda tramite l'algoritmo di prestito. \\
Inoltre, per capire l'effetto del traffic shaping nella rete, è utile guardare anche i risultati senza l'utilizzo di HTB.
\begin{center}
  \begin{tabular}{ |c|c|c| }
    \hline
    \textbf{Classe} & \textbf{Tempo} \\
    \hline
    Sito Web & $22.81s$ \\
    \hline
    Streaming & $23.04s$ \\
    \hline
    File Hosting & $24.64s$ \\
    \hline
  \end{tabular}
\end{center}
Si può notare come i tempi, simili tra loro, rappresentino un terzo della banda totale, ossia circa 33Mbit/s per ciascun nodo, 
dato che $\frac{100 \text{MB} \times 8}{33 \text{Mbit}/s} = 24s$. \\
\vspace*{0.01\textheight}
In conclusione, applicando una disciplina di accodamento come HTB per fare traffic shaping,
ha permesso al servizio del sito web di avere a disposizione una banda maggiore, terminando il trasferimento per primo. 
Invece, ha rallentato il servizio di file hosting, dedicandogli una banda minore. \\
L'utilizzo di Kathará per creare la rete, ha permesso di calcolare i tempi reali in una LAN nelle diverse modalità proposte.
\begin{center}
  \begin{tabular}{ |c|c|c|c| }
    \hline
    \textbf{Servizio} & \textbf{Tempo con HTB} & \textbf{Tempo senza HTB}\\
    \hline
    Sito Web & $19.20s$ & $22.81s$ \\
    \hline
    Streaming & $23.33s$ & $23.04s$ \\
    \hline
    File Hosting & $48.47s$ & $24.64s$ \\
    \hline
  \end{tabular}
\end{center}

\newpage
\section{Dipendenza dai moduli del kernel host}
L'uso di container Docker rende la simulazione leggera e flessibile, permettendo di scegliere per ogni nodo un'immagine diversa, ma crea un limite imporntante.
A differenza delle macchine virtuali, i container non eseguono un kernel proprio, ma condividono quello dell'host su cui gira Docker, 
ciò che un nodo può fare a livello di rete dipende dai moduli caricati nel kernel dell'host. \\
\vspace*{0.01\textheight}
Il comando tc, eseguito dentro un container, iniva solamente richieste al kernel, non implementa effettivamente le discipline di accodamento.
Le qdisc, le classi e i filtri sono realizzati da diversi moduli del kernel:
\begin{center}
  \begin{tabular}{ |c|c| }
    \hline
    \textbf{Modulo} & \textbf{Descrizione} \\
    \hline
    \texttt{sch\_tbf} & Modulo per qdisc TBF \\
    \hline
    \texttt{sch\_htb} & Modulo per qdisc HTB \\
    \hline
    \texttt{cls\_u32} & Modulo per filtri \texttt{u32} \\
    \hline
  \end{tabular}
\end{center}



